\scanpage{41}
% archive sheet number: 40
les él. de
\[
E^-_{\deleted{\unsure{f}}}(i)_{f_i}/(\sigma_0,\ldots,\widehat{\sigma_j},\ldots,\sigma_{i-1}),\quad\text{i.e.}
\]
les él. de
\[
E^-(i)/(\sigma_0,\ldots,\widehat{\sigma_j},\ldots,\sigma_{i-1})
\simeq E/(\sigma_0,\ldots,\widehat{\sigma_j},\ldots,\widehat{\sigma_i},\ldots,\sigma_n)
\]
qui \deleted{\unsure{au delà}} s'envoient dans $f_i\inD_i$, i.e. les drapeaux de types $(j,i)$ dont la composante de \deleted{degré} dim $i$ est $f_i$. Plus gén, ses drapeaux de type $(j_1<\cdots<j_p)$ (i.e. $0\leq j_1,\deleted{\unsure{j}}j_p\leq\unsure{i-1|i}$) sont les drapeaux de type $(j_1,\ldots,j_p,i)$ de $\Pi$ \deleted{dont} la composante de degré $i$ est $f_i$. On l'appelle \unsure{le|la} \inserted{\uline{carte}} ``\uline{bord}'' de la facette $f_i$, \unsure{notée|noté} $\partial f_i$.

\unsure{Dualement|Duellement}, considérons
\[
E^+(i)=E/G_n^-(i)\inserted{=D_{i,i+1,\ldots,n}}\longrightarrow D_i
\]
sur lequel $G_n^+(i)\simeq\deleted{\unsure{n-i}}G_{n-i-1}$ opère, le quotient étant $D_i$. On trouve une moquette \inserted{de $(n-i-1)$-cartes} (sauf erreur), dont les facettes de dim $j$ $(0\leq j\leq n-i-1)$ ``sont'' les drapeaux de type $(i,i+j+1)$ \inserted{de $\Pi$} \deleted{dont} la composante de dim $i$ est $f_i$.

On peut aussi, pour \deleted{deux facettes} \inserted{un drapeau \inserted{$d_{ij}$} de type $(i,j)$, \unsure{défini par|définissant}}
\[
f_i<f'_j\qquad(\deleted{i<j}\inserted{j\geq i+2})
\]
considérer
\[
\deleted{\unsure{E(i,j)=E/(\sigma_0-\sigma_{i-1};\sigma_{j+1},\ldots,\sigma)}}
\]
